\documentclass[a4paper,12pt]{article}
\usepackage[utf8]{inputenc}
\usepackage[T1]{fontenc}
\usepackage{lmodern}
\usepackage{amsmath}
\usepackage{geometry}
\usepackage{booktabs}
\usepackage{graphicx}
\usepackage{enumitem}
\usepackage{xcolor}
\geometry{margin=2.5cm}
\setlength{\parindent}{0pt}
\setlength{\parskip}{0.6em}

\title{IoT - Challenge 2 - Part 2}
\author{Raul Agolli - Team Leader \\\\ Gabriele Lorenzo Singe}
\date{April 2026}

\begin{document}

\maketitle

\section{Introduction}

This report presents the solution to Part 2 of Challenge 2.
The system consists of:
\begin{itemize}[leftmargin=1.5cm]
    \item a battery-powered temperature sensor transmitting one measurement every 2 minutes,
    \item a battery-powered actuator (fan controller) receiving temperature data and deciding whether to switch ON/OFF,
    \item a mains-powered gateway supporting only MQTT.
\end{itemize}

The communication can take place either:
\begin{itemize}[leftmargin=1.5cm]
    \item directly using CoAP,
    \item through the gateway using MQTT.
\end{itemize}

Communication parameters:
\begin{itemize}
    \item Payload size: 12 bytes
    \item Topic/resource length: 8 bytes
\end{itemize}



The given message sizes are:
\begin{center}
\begin{tabular}{l c}
\toprule
\textbf{CoAP message} & \textbf{Size [B]} \\
\midrule
CON PUT Request & 70 \\
PUT ACK & 15 \\
\bottomrule
\end{tabular}
\hspace{1cm}
\begin{tabular}{l c}
\toprule
\textbf{MQTT message} & \textbf{Size [B]} \\
\midrule
Subscribe & 65 \\
SubAck & 55 \\
Publish & 75 \\
PubAck & 50 \\
Connect/Disconnect & 60 \\
ConnectAck & 50 \\
\bottomrule
\end{tabular}
\end{center}

The energy parameters are:
\[
E_{TX} = 50 \text{ nJ/bit}, \qquad E_{RX} = 58 \text{ nJ/bit}
\]

The assumptions used are:
\begin{enumerate}[leftmargin=1.5cm]
    \item CoAP uses Confirmable (CON) messages.
    \item MQTT uses QoS 1.
    \item Communication is ideal, with no losses.
    \item For EQ1, nodes establish the connection once and remain connected.
    \item Only the energy consumed by sensor and actuator is considered.
\end{enumerate}

Since the sensor transmits one measurement every 2 minutes, in one hour the total number of measurements is:
\[
N = \frac{60}{2} = 30
\]

\section{EQ1 --- Communication energy over 1 hour}

\subsection{EQ1a --- CoAP direct communication}

For each measurement:
\begin{itemize}[leftmargin=1.5cm]
    \item the sensor transmits one CON PUT request of 70 B,
    \item the actuator receives the 70 B request,
    \item the actuator transmits one PUT ACK of 15 B,
    \item the sensor receives the 15 B ACK.
\end{itemize}

Therefore, the energy consumed for one CoAP exchange is:
\[
E_{\text{CoAP},1} =
70 \cdot 8 \cdot E_{TX}
+ 70 \cdot 8 \cdot E_{RX}
+ 15 \cdot 8 \cdot E_{TX}
+ 15 \cdot 8 \cdot E_{RX}
\]

\[
E_{\text{CoAP},1} = 73.44 \ \mu J
\]

Since there are 30 measurements in one hour:
\[
E_{\text{CoAP},1h} = 30 \cdot 73.44 = 2203.2 \ \mu J
\]

\[
\boxed{EQ1a = 2203.2 \ \mu J}
\]

\subsection{EQ1b --- MQTT communication via gateway}

For EQ1, nodes establish the MQTT connection once and remain connected for the whole hour.

\subsubsection*{Sensor energy}

At the beginning, the sensor:
\begin{itemize}[leftmargin=1.5cm]
    \item transmits one CONNECT message of 60 B,
    \item receives one CONNACK message of 50 B.
\end{itemize}

Then, for each of the 30 measurements:
\begin{itemize}[leftmargin=1.5cm]
    \item it transmits one PUBLISH message of 75 B,
    \item it receives one PUBACK message of 50 B.
\end{itemize}

Thus:
\[
E_{\text{sens}} =
(60 + 30 \cdot 75)\cdot 8 \cdot E_{TX}
+
(50 + 30 \cdot 50)\cdot 8 \cdot E_{RX}
\]

\[
E_{\text{sens}} = 1643.2 \ \mu J
\]

\subsubsection*{Actuator energy}

At the beginning, the actuator:
\begin{itemize}[leftmargin=1.5cm]
    \item transmits one CONNECT message of 60 B,
    \item receives one CONNACK message of 50 B,
    \item transmits one SUBSCRIBE message of 65 B,
    \item receives one SUBACK message of 55 B.
\end{itemize}

Then, for each of the 30 measurements:
\begin{itemize}[leftmargin=1.5cm]
    \item it receives one PUBLISH message of 75 B,
    \item it transmits one PUBACK message of 50 B.
\end{itemize}

Therefore:
\[
E_{\text{act}} =
(60 + 65 + 30 \cdot 50)\cdot 8 \cdot E_{TX}
+
(50 + 55 + 30 \cdot 75)\cdot 8 \cdot E_{RX}
\]

\[
E_{\text{act}} = 1742.72 \ \mu J
\]

\subsubsection*{Total MQTT energy}

The total energy over one hour is:
\[
E_{\text{MQTT},1h} = E_{\text{sens}} + E_{\text{act}}
\]

\[
\boxed{EQ1b = 3385.92 \ \mu J}
\]

\section{EQ2 --- Protocol choice with actuator waking up every 30 minutes}

In this scenario, the actuator wakes up only once every 30 minutes, while the sensor still generates a temperature measurement every 2 minutes.
The goal is to minimize the long-term communication energy consumption of the battery-powered devices.

In this case, MQTT is the preferable solution, because the broker can store the latest published value and deliver it as soon as the actuator reconnects.
This avoids forcing the actuator to wake up every 2 minutes or to miss the most recent information.

A suitable configuration is:
\begin{itemize}[leftmargin=1.5cm]
    \item MQTT with QoS 1,
    \item retained publish,
    \item persistent session (\texttt{clean session = false}).
\end{itemize}

With this choice, when the actuator wakes up every 30 minutes, it can immediately obtain the latest available temperature value while minimizing its active communication time.

\bigskip
\noindent
\textbf{Answer for EQ2:}

\begin{quote}
We would choose MQTT via the gateway, with QoS 1, retained publish, and persistent session (\texttt{clean session = false}). The optimized metric is actuator communication energy: when it wakes up every 30 minutes, it can immediately receive the latest temperature value.
\end{quote}

\end{document}